\documentclass[aspectratio=169]{beamer}
\usetheme{Madrid}
\usecolortheme{default}
\usepackage[utf8]{inputenc}
\usepackage[T5]{fontenc}
\usepackage{graphicx}
\usepackage{hyperref}
\usepackage{booktabs}
\usepackage{amsmath}

\setbeamertemplate{caption}[numbered]
\renewcommand{\figurename}{Hình}
\renewcommand{\thefigure}{9.\arabic{figure}}

\title[Chương 9: Kiến trúc ConvNet]{Học Sâu (Deep Learning) \\ \vspace{0.5cm} \Large Chương 9: Các Mẫu Kiến trúc Mạng Tích chập \\ (ConvNet Architecture Patterns)}
\author{Đại học Đông Á}
\date{\today}

\begin{document}

\begin{frame}
    \titlepage
\end{frame}

\begin{frame}{Nội dung chương}
    \tableofcontents
\end{frame}

\section{1. Nguyên lý Thiết kế Kiến trúc}

\begin{frame}{Không gian Giả thuyết (Hypothesis Space)}
    \begin{itemize}
        \item \textbf{Kiến trúc Mô hình (Model Architecture)} là nghệ thuật lắp ráp các tầng mạng (Layers).
        \item Cấu trúc này định nghĩa \textbf{Không gian Giả thuyết} (Hypothesis Space). Thuật toán Gradient Descent sẽ tìm kiếm bộ tham số (Weights) tối ưu bên trong không gian này.
        \item Một kiến trúc tốt giống như một bản thiết kế thông minh, giúp đóng gói \textbf{Tri thức tiên nghiệm (Prior Knowledge)} của kỹ sư vào mô hình, làm thu hẹp không gian tìm kiếm và tăng tốc độ hội tụ.
    \end{itemize}
\end{frame}

\begin{frame}{Công thức MHR: Modularity, Hierarchy, Reuse}
    \begin{columns}
        \column{0.5\textwidth}
        Để quản lý một hệ thống phức tạp, ta dùng nguyên lý \textbf{MHR}:
        \begin{itemize}
            \item \textbf{Modularity (Mô-đun hóa):} Đóng gói các lớp thành các khối (Blocks).
            \item \textbf{Hierarchy (Phân cấp):} Xếp chồng các khối thành tháp (Pyramid).
            \item \textbf{Reuse (Tái sử dụng):} Tái sử dụng cùng một cấu trúc khối nhiều lần.
        \end{itemize}
        
        \column{0.5\textwidth}
        \begin{figure}
            \centering
            \includegraphics[width=\textwidth,height=0.6\textheight,keepaspectratio]{../../images/ch09/complex_systems.e89aaf78.png}
            \caption{Hệ thống phức tạp sinh học tuân theo nguyên lý MHR}
        \end{figure}
    \end{columns}
\end{frame}

\begin{frame}{Tháp Đặc trưng trong Mô hình Xception}
    \begin{columns}
        \column{0.5\textwidth}
        \begin{itemize}
            \item Các mạng CNN hiện đại (như Xception, ResNet) tuân thủ chặt chẽ nguyên lý MHR.
            \item Cấu trúc Lặp lại: \texttt{SeparableConv} $\rightarrow$ \texttt{SeparableConv} $\rightarrow$ \texttt{MaxPooling}.
            \item \textbf{Tháp không gian:} Kích thước bản đồ (Feature Maps) giảm dần ($299 \times 299 \rightarrow 19 \times 19$).
            \item \textbf{Tháp chiều sâu:} Số lượng bộ lọc (Filters) tăng dần ($32 \rightarrow 728$).
        \end{itemize}
        
        \column{0.5\textwidth}
        \begin{figure}
            \centering
            \includegraphics[width=\textwidth,height=0.7\textheight,keepaspectratio]{../../images/ch09/xception_entry_flow_pyramid.4701a5de.png}
            \caption{Sự phân cấp của Khối đầu vào Xception}
        \end{figure}
    \end{columns}
\end{frame}

\section{2. Batch Normalization (Chuẩn hóa Hàng loạt)}

\begin{frame}{Vấn đề: Internal Covariate Shift}
    \begin{itemize}
        \item Trong một mạng nhiều tầng, giá trị kích hoạt (Activations) ở một tầng phụ thuộc vào tầng trước nó.
        \item Khi trọng số của tầng đầu thay đổi trong quá trình huấn luyện (Gradient Descent), \textbf{phân phối dữ liệu} ở tầng sau cũng bị xô lệch theo.
        \item Hiện tượng này gọi là \textbf{Internal Covariate Shift}, làm giảm tốc độ hội tụ và khiến mạng sâu rất khó huấn luyện.
        \item Giải pháp: Chuẩn hóa dữ liệu ngay trong quá trình luân chuyển giữa các tầng.
    \end{itemize}
\end{frame}

\begin{frame}{Toán học của Batch Normalization}
    Lớp \texttt{BatchNormalization} chuẩn hóa từng Batch dữ liệu trong lúc huấn luyện:
    \begin{enumerate}
        \item Tính Kỳ vọng ($\mu$) của Batch: $\mu = \frac{1}{m} \sum_{i=1}^m x_i$
        \item Tính Phương sai ($\sigma^2$) của Batch: $\sigma^2 = \frac{1}{m} \sum_{i=1}^m (x_i - \mu)^2$
        \item Chuẩn hóa: $\hat{x}_i = \frac{x_i - \mu}{\sqrt{\sigma^2 + \epsilon}}$
        \item Chuyển tỷ lệ và dịch chuyển (Scale \& Shift): $y_i = \gamma \hat{x}_i + \beta$
    \end{enumerate}
    Trong đó $\gamma$ và $\beta$ là hai tham số học được trong quá trình huấn luyện, đảm bảo mô hình có thể khôi phục lại bất kỳ phân phối nào nếu cần thiết.
\end{frame}

\begin{frame}{Ứng dụng Batch Normalization trong Keras}
    \begin{itemize}
        \item Vị trí lý tưởng: Thông thường đặt ngay sau tầng \texttt{Conv2D} hoặc \texttt{Dense} và trước khi qua Hàm kích hoạt (Activation).
        \item \textbf{Cơ chế hoạt động kép:}
        \begin{itemize}
            \item Khi Huấn luyện (\texttt{training=True}): Dùng $\mu$ và $\sigma^2$ của Batch hiện tại. Đồng thời cập nhật trung bình động trượt (Moving Average) cho toàn bộ tập dữ liệu.
            \item Khi Suy luận (\texttt{training=False}): Dùng giá trị trung bình động (Moving Average) đã tích lũy thay vì tính toán lại, để bảo đảm tính ổn định tuyệt đối cho đầu ra.
        \end{itemize}
    \end{itemize}
\end{frame}

\section{3. Kết nối Thặng dư (Residual Connections)}

\begin{frame}{Vấn đề: Suy biến Đạo hàm (Gradient Vanishing)}
    \begin{itemize}
        \item Nếu bạn liên tục thêm hàng chục tầng Convolutional vào một mô hình (Deep Network), Độ chính xác huấn luyện không tăng mà lại \textbf{giảm xuống}.
        \item Lý do: Thuật toán Backpropagation truyền đạo hàm ngược từ Hàm Loss về lại các tầng đầu tiên thông qua Quy tắc Chuỗi (Chain Rule).
        \item Phép nhân liên tiếp nhiều số nhỏ hơn 1 khiến đạo hàm giảm mạnh dần về số 0. Mạng bị "bất động", không thể học được.
    \end{itemize}
\end{frame}

\begin{frame}{Giải pháp Kết nối Thặng dư (Residual Connections)}
    \begin{columns}
        \column{0.5\textwidth}
        \begin{itemize}
            \item Phát minh của Kaiming He (2015 - Đạt giải CVPR Best Paper).
            \item Cung cấp một \textbf{đường tắt (Shortcut)} cho thông tin nhảy vọt qua một (hoặc nhiều) lớp không cần tính toán.
            \item Đạo hàm lúc này sẽ chảy mượt mà qua các đường tắt này về tận đầu mạng mà không bị hao hụt.
        \end{itemize}
        
        \column{0.5\textwidth}
        \begin{figure}
            \centering
            \includegraphics[width=\textwidth,height=0.7\textheight,keepaspectratio]{../../images/ch09/residual_connection.0524fdc4.png}
            \caption{Cơ chế hoạt động của Khối Residual}
        \end{figure}
    \end{columns}
\end{frame}

\begin{frame}{Toán học của Khối Residual (ResNet)}
    Giả sử $x$ là đầu vào của khối, và $F(x)$ là kết quả tính toán qua 2 lớp Convolution.
    \begin{itemize}
        \item Mạng truyền thống (Plain Network): Đầu ra là $y = F(x)$
        \item Mạng Thặng dư (Residual Network): Đầu ra là $y = F(x) + x$
    \end{itemize}
    Thay vì học ánh xạ tổng $y$, khối Convolution $F(x)$ lúc này chỉ cần học phần \textbf{sai khác (thặng dư - Residual)} giữa đầu ra và đầu vào. 
    Việc học "sự thay đổi nhỏ" luôn dễ dàng hơn học "mọi thứ từ đầu".
\end{frame}

\begin{frame}{Giải quyết Bất đồng bộ Kích thước}
    Khi áp dụng \texttt{MaxPooling} hoặc thay đổi Số lượng Bộ lọc, kích thước của $F(x)$ và $x$ sẽ khác nhau, không thể áp dụng phép cộng trực tiếp.
    \begin{itemize}
        \item Giải pháp: Sử dụng \textbf{Linear Projection} (Tích chập $1 \times 1$) không kích hoạt (No Activation).
        \item $y = F(x) + W_s \cdot x$
        \item Keras: \texttt{layers.Conv2D(filters, kernel\_size=1, strides=2, padding="same")(x)}
    \end{itemize}
\end{frame}

\section{4. Tích chập Tách rời (Separable Convolution)}

\begin{frame}{Vấn đề của Tích chập Truyền thống (Conv2D)}
    \begin{itemize}
        \item Tích chập chuẩn thực hiện đồng thời 2 việc trong cùng 1 bước:
        \begin{enumerate}
            \item Học mẫu không gian (Spatial Pattern) trên mặt phẳng 2D.
            \item Tổ hợp thông tin từ nhiều Kênh khác nhau (Cross-channel).
        \end{enumerate}
        \item Sự gắn kết này dẫn đến chi phí tính toán cực kỳ lớn, với số lượng tham số là: $K^2 \times C_{in} \times C_{out}$.
        \item Liệu Không gian 2D (Spatial) và Chiều sâu Kênh (Channel) có thực sự phụ thuộc chặt chẽ như vậy?
    \end{itemize}
\end{frame}

\begin{frame}{Khái niệm Separable Convolution}
    \begin{columns}
        \column{0.5\textwidth}
        Tách rời phép tính thành 2 pha độc lập:
        \begin{enumerate}
            \item \textbf{Depthwise Convolution:} Quét từng Kênh một cách riêng biệt (Không giao thoa Kênh).
            \item \textbf{Pointwise Convolution:} Dùng bộ lọc $1 \times 1$ để nhào trộn và tổng hợp các kênh.
        \end{enumerate}
        Giả định: Đặc trưng Không gian (Mắt, Mũi) và Đặc trưng Kênh (Đỏ, Xanh) có tính độc lập cao.
        
        \column{0.5\textwidth}
        \begin{figure}
            \centering
            \includegraphics[width=\textwidth,height=0.7\textheight,keepaspectratio]{../../images/ch09/depthwise_separable_conv.5d1929bd.png}
            \caption{Tích chập Tách rời chiều sâu}
        \end{figure}
    \end{columns}
\end{frame}

\begin{frame}{Ưu điểm vượt trội về Toán học và Tốc độ}
    \begin{itemize}
        \item Bằng cách chia rẽ quá trình tính toán, số lượng Tham số (Parameters) được thu gọn:
        \item Tích chập chuẩn: $K^2 \times C_{in} \times C_{out}$
        \item Tích chập Tách rời: $K^2 \times C_{in} + C_{in} \times C_{out}$
        \item \textbf{Ví dụ:} Kernel $3 \times 3$, $C_{in}=64$, $C_{out}=64$.
        \begin{itemize}
            \item Lớp \texttt{Conv2D} cần: $3^2 \times 64 \times 64 = 36,864$ tham số.
            \item Lớp \texttt{SeparableConv2D} chỉ cần: $(3^2 \times 64) + (64 \times 64) = 576 + 4,096 = 4,672$ tham số.
        \end{itemize}
        \item $\rightarrow$ Tiết kiệm gần 90\% năng lực tính toán nhưng vẫn duy trì độ chính xác.
    \end{itemize}
\end{frame}

\section{5. Thực hành: Xây dựng Mini-Xception}

\begin{frame}{Mô hình Mini-Xception cho Chó \& Mèo}
    Lắp ráp 3 bảo bối của Học sâu vào chung một mạng (Mini-Xception):
    \begin{itemize}
        \item \textbf{Lõi xử lý:} Thay vì dùng \texttt{Conv2D}, ta sử dụng toàn bộ \texttt{SeparableConv2D} để giảm chi phí tính toán.
        \item \textbf{Bình ổn hóa:} Mỗi bước tích chập đều kèm theo \texttt{BatchNormalization} để ổn định tốc độ học.
        \item \textbf{Thông suốt Đạo hàm:} Mỗi Khối (Block) được bắc cầu bởi một \textbf{Residual Connection} sử dụng Projection $1 \times 1$.
    \end{itemize}
\end{frame}

\begin{frame}{Đánh giá Độ chính xác (Accuracy)}
    \begin{columns}
        \column{0.5\textwidth}
        \begin{itemize}
            \item Nhờ cấu trúc siêu ưu việt của \texttt{SeparableConv} và Residual, mô hình tự thiết kế của chúng ta đạt \textbf{Accuracy lên đến 90\%} (Chỉ huấn luyện lại từ đầu với Data Augmentation, KHÔNG DÙNG Transfer Learning).
            \item Con số này đánh bại hoàn toàn mô hình Basic CNN (82\%) ở Chương 8.
        \end{itemize}
        
        \column{0.5\textwidth}
        \begin{figure}
            \centering
            \includegraphics[width=\textwidth,height=0.5\textheight,keepaspectratio]{../../images/ch09/training-and-validation-xception-acc.967bf65c.png}
            \caption{Accuracy của Mini-Xception}
        \end{figure}
    \end{columns}
\end{frame}

\begin{frame}{Đánh giá Mất mát (Loss)}
    \begin{columns}
        \column{0.5\textwidth}
        \begin{itemize}
            \item Biểu đồ Loss cho thấy đường Validation Loss đi xuống ổn định trong suốt quá trình.
            \item \texttt{BatchNormalization} giúp thuật toán hội tụ cực nhanh.
            \item Thiết kế này là tiêu chuẩn công nghiệp (Industry Standard) hiện nay trên Keras cho hầu hết mọi tác vụ Thị giác máy tính.
        \end{itemize}
        
        \column{0.5\textwidth}
        \begin{figure}
            \centering
            \includegraphics[width=\textwidth,height=0.5\textheight,keepaspectratio]{../../images/ch09/training-and-validation-xception-loss.bc0f5982.png}
            \caption{Loss của Mini-Xception}
        \end{figure}
    \end{columns}
\end{frame}

\section{Tổng kết}

\begin{frame}{Tổng kết Chương 9}
    Để xây dựng một mô hình CNN xuất sắc từ con số 0, hãy tuân theo các nguyên tắc sau:
    \begin{itemize}
        \item Tổ chức Mô hình theo dạng Tháp Đặc trưng (Feature Pyramid). Lặp lại Khối (Blocks) thay vì tự nối từng Layer ngẫu nhiên (Quy tắc MHR).
        \item Sử dụng \textbf{Separable Convolution} thay cho Convolution tiêu chuẩn (Tính toán cực nhẹ).
        \item Kết hợp \textbf{Batch Normalization} để hệ thống hội tụ mượt mà không vấp ngã.
        \item Tạo các \textbf{Kết nối Thặng dư (Residual Connections)} để dòng chảy Gradient truyền xa xuyên suốt mạng sâu 100 tầng.
    \end{itemize}
\end{frame}

\end{document}
